problem:  
Let $E_k \to S^4$ be the smooth $S^3$–bundle associated to the principal $SU(2)$–bundle over $S^4$ with second Chern number $c_2 = k \in \mathbb Z$.  
For a given $SU(2)$–connection $A$, let $g_{\lambda,A}$ denote the standard connection metric on $E_k$ with fiber radius $\lambda$.  

Now consider the **Cheeger deformation** of each $g_{\lambda,A}$ along the canonical fiberwise $SU(2)$–action, producing a family of metrics $g_{\lambda,A,t}$ on $E_k$, where $t>0$ is the deformation parameter.  For homogenous and biquotient cases (certain small $|k|$), positive sectional curvature metrics are known at special values of $(\lambda,t)$.  

**Question:**  
For general $k\in\mathbb Z$, does there exist a connection $A$ and parameters $\lambda,t>0$ such that the Cheeger–deformed metric $g_{\lambda,A,t}$ on $E_k$ has strictly positive sectional curvature on all $2$‑planes?  
Equivalently, can the Cheeger deformation destroy all remnants of negative mixed curvature terms arising in the O’Neill formulas for nonfat connections, or does the topological twisting class $k$ impose an intrinsic obstruction to positivity even under all such deformations?  

Why is it a "good" problem:  
This formulation is **precise, geometric, and topologically specific.**  It concretely examines whether the Cheeger deformation—a canonical curvature‑enhancing process—can eliminate the negative curvature modes inherent in nonfat connection metrics on a distinguished and nontrivial family of sphere bundles ($S^3$–bundles over $S^4$).  The class parameter $k$ provides an explicit topological marker linking differential geometry to algebraic topology.

The problem strikes directly at a fault line in current knowledge: while Cheeger deformation is known to improve certain curvature components, its ultimate power to force full positivity on nonhomogeneous fiber bundles remains unresolved.  A positive solution would reveal a methodical route from bundle topology to positively curved metrics through explicit geometric flow; a negative one would pinpoint a topological curvature obstruction tied to the Chern number $k$.

Its statement is simple yet profound—identifying just one familiar geometric operation (Cheeger deformation) and one numerical invariant ($k$)—but its implications would bridge Riemannian submersion curvature, gauge fields on $S^4$, and the global topology of positively curved manifolds.  This blend of conceptual clarity, cross‑field resonance, and unresolved depth makes it a genuinely “good” research‑level problem.